\section{Architecture}

\subsection{Baseline Gated FFN}

A gated Transformer FFN \cite{shazeer2020gptj} maps hidden states
$X\in\R^{B\times S\times d}$ to
\begin{equation}
Y = \left(\phi(XW_{\mathrm{gate}})\odot XW_{\mathrm{up}}\right)W_{\mathrm{down}},
\end{equation}
where $W_{\mathrm{up}}, W_{\mathrm{gate}}\in\R^{d\times m}$,
$W_{\mathrm{down}}\in\R^{m\times d}$, and $\phi$ is typically SiLU. With fp16
weights, these three matrices store exactly $6dm$ bytes, excluding metadata and
alignment.

\subsection{Ephemeral-Core FFN}

For $p\in\{u,g,d\}$, let $U_p\in\R^{d_{\rm in}\times r}$,
$V_p\in\R^{r\times d_{\rm out}}$ be persistent factors, and a generator
emits $M_{p,\mathcal G}\in\mathcal M_r$ from causal group context, layer id,
and projection id:
\begin{equation}
\widetilde Y_{\mathcal G}=
\bigl(\phi(\Pi_{g,\mathcal G}(X_{\mathcal G}))\odot
\Pi_{u,\mathcal G}(X_{\mathcal G})\bigr)
U_dM_{d,\mathcal G}V_d.
\label{eq:core-ffn}
\end{equation}
The structured family controls descriptor size and expressivity: a diagonal
core $\operatorname{Diag}(z)$ emits $r$ scalars per projection with
$\Theta(r)$ core work per token; block-diagonal ($b$-blocks, $rb$ scalars),
low-rank-plus-identity ($I+AB^\top$, $2rk$), and dense ($r^2$) families trade
expressivity for payload. We adopt the diagonal family as canonical.
Generating complete factors $(U_{p,\mathcal G},V_{p,\mathcal G})$ is a more
ambitious alternative, not the implemented method: it emits $r(d+m)$ scalars
per projection and exposes a severe reuse requirement
(Section~\ref{sec:analysis})---which is exactly why a small core, rather
than a full generated factor pair, is the practical object for the proposed
lifetime contract.

\subsection{Causal Token-Local Replacement}
\label{sec:causal-token}

The historical \texttt{expert\_residual} mode formed its context as the
full-sequence-and-batch mean $\bar x$, which is not valid for causal
evaluation (position $t$'s control depends on later positions and unrelated
examples); we retain it only for checkpoint compatibility. The replacement
mode \texttt{token\_diag\_core} applies a controller independently per
pre-FFN token state:
\begin{equation}
(z_u,z_g,z_d)_{b,t}=G_\theta(x_{b,t}),\qquad
P(x;U,V,z)=((xU)\odot z)V,
\label{eq:diag-core}
\end{equation}
using the bounded parameterization $z=\mathbf 1+a\tanh(h)$, with $a$ a
recorded amplitude. It emits $3r$ values (vs.\ $6r+K$), uses no residual
expert bank, and has $q=1$: because a decoder state depends only on its
prefix, the output is independent of later positions and companion batch
examples. A future shared core with $q>1$ is causal only if its context is
available before the first token in the group. Tests compare full-sequence
and prefix-only calls, perturb unrelated batch members, verify
Eq.~\eqref{eq:diag-core} explicitly, and round-trip checkpoints. The legacy
mode remains loadable only for compatibility; its head cannot generally be
converted into an affine $3r$ head, so the redesigned mode needs fresh
training (redundancy proposition, Appendix~\ref{app:causal}).

\subsection{Descriptor Lifecycle}

The inference operator enforces the following contract:
(i) a causal router dispatches each call before either FFN branch launches
(the unselected branch must not run); (ii) persistent $U_p,V_p,\theta$ may
traverse HBM and cache, while the generator emits $M_{p,\mathcal G}$ directly
into kernel-local scratch; (iii) the fused projection consumes $M_{p,\mathcal G}$
without returning, aliasing, or registering it as model state, and without
storing it in a global-memory descriptor buffer; and (iv) scratch is
overwritten before the next group, with device counters and an instrumented
address range distinguishing this arm from an otherwise identical
materialized-core arm.

The contract is made operational by a single entry point
$\mathsf{ephemeral\_ffn}$ with signature
\begin{equation}
y_{\mathcal G},\ \mathsf{counters}\;\leftarrow\;
\mathsf{ephemeral\_ffn}\!\left(X_{\mathcal G};\ U_p,V_p,\theta;\ g(x);\ \mathsf{mode}\right),
\label{eq:operator-api}
\end{equation}
where $g(x)\in\{0,1\}$ is the causal dispatch decision and $\mathsf{mode}$
selects inference, save-core, or rematerialize-core backward semantics. Its
preconditions are that $g(x)$ is computed before either branch launches, that
$U_p,V_p,\theta$ are durable model state, and that the core family
$\mathcal M_r$ and output tolerance are declared with the operator. Its
postconditions are that the returned $y_{\mathcal G}$ matches the selected
branch within tolerance, that in inference mode no core buffer is returned,
aliased, registered as model state, or written outside the declared scratch
range, and that $\mathsf{counters}$ reports reads/writes of the core region.
These last postconditions are the testable content of the no-store contract: a
kernel that fails them is an ordinary materialized generator, not an \mwg
operator. The target memory hierarchy is depicted in
Figure~\ref{fig:lifecycle} (Appendix~\ref{app:lifecycle}).

Training requires different semantics.  \texttt{save-core} retains each core
until backward and weakens the strict no-store claim; alternatively,
\texttt{rematerialize-core} retains the causal source activation and
deterministic RNG/replay metadata, regenerates the core in backward, and
propagates generator gradients through the replayed computation. Neither
option is implemented by the target kernel, so we restrict the principal
lifecycle claim to inference.